\documentclass[a4paper,12pt]{article}
\usepackage[T2A]{fontenc}
\usepackage[utf8]{inputenc}
\usepackage[russian]{babel}
\usepackage{amsmath,amssymb,mathtools}
\usepackage{icomma}
\usepackage{geometry}
\geometry{left=21mm,right=21mm,top=21mm,bottom=22mm,headheight=26pt,headsep=4.5mm}
\usepackage{microtype}
\usepackage{tikz}
\usetikzlibrary{calc,arrows.meta,positioning,shapes.geometric}
\usepackage{tkz-euclide}
\usepackage{tabularx,booktabs,array}
\usepackage{enumitem}
\usepackage{fancyhdr}
\usepackage{setspace}
\usepackage{parskip}
\usepackage{needspace}
\usepackage{xcolor}
\usepackage[hidelinks]{hyperref}
\definecolor{accent}{HTML}{007F87}
\definecolor{darkaccent}{HTML}{00646B}
\definecolor{textgray}{HTML}{363636}
\definecolor{lightgray}{HTML}{657178}
\setstretch{1.20}
\setlength{\parindent}{0pt}
\setlength{\parskip}{6.4pt plus 1pt minus 1pt}
\setlist[enumerate]{leftmargin=*,itemsep=7pt,topsep=4pt}
\setlist[itemize]{leftmargin=*,itemsep=2pt,topsep=3pt}
\pagestyle{fancy}
\fancyhf{}
\fancyhead[L]{\small\color{textgray}Призмы: объём и отношения}
\fancyhead[R]{\small\color{textgray}Ксения Клыкова}
\fancyfoot[C]{\small\thepage}
\renewcommand{\headrulewidth}{0.45pt}
\renewcommand{\headrule}{\hbox to\headwidth{\color{accent}\leaders\hrule height \headrulewidth\hfill}}
\newcommand{\topic}[1]{\needspace{6\baselineskip}\section{#1}\vspace{-1.5mm}\noindent\textcolor{accent}{\rule{\linewidth}{0.45pt}}\vspace{-2mm}}
\newcommand{\key}[1]{\textcolor{darkaccent}{\textbf{#1}}}
\newcommand{\ex}[1]{\par\medskip\noindent\textcolor{darkaccent}{\textbf{#1}}\par\smallskip}
\newcommand{\mat}[1]{\[#1\]}
\begin{document}
\thispagestyle{empty}
\begin{center}
\vspace*{20mm}
{\large\color{darkaccent} МАТЕМАТИКА}\par
\vspace{20mm}
{\LARGE\color{textgray} Чек-лист}\par
\vspace{15mm}
{\LARGE\color{accent} Призмы: объём, боковая\\[2mm] поверхность и отношения}\par
\vspace{17mm}
{\normalsize Подготовка к ЕГЭ по профильной математике}\par
\vfill
{\large Ксения Клыкова}\par
\vspace{15mm}
{\normalsize 24 ноября 2026 г.}\par
\vspace{16mm}
{\small Лёвин Артём Александрович\\[1.5mm]эксклюзивно для Ксении Клыковой}\par
\vspace{10mm}
\begin{tikzpicture}
\draw[accent,line width=0.85pt] (0,0) .. controls (1.4,0.66) and (2.4,-0.35) .. (3.95,0.04)
 .. controls (5.2,0.64) and (6.5,-0.4) .. (8.1,0.08);
\end{tikzpicture}
\vspace*{3mm}
\end{center}
\clearpage

\topic{Базовая формула объёма призмы}
Для \emph{любой} призмы объём равен произведению площади основания на высоту:
\[
\boxed{V=S_{\text{осн}}\,h}.
\]
Высота $h$ --- перпендикулярное расстояние между плоскостями оснований. У прямой призмы она равна боковому ребру. У наклонной призмы эти длины могут различаться.

\key{Что нужно знать об основании:}
\begin{itemize}
\item правильный треугольник со стороной $a$: \(\displaystyle S_{\triangle}=\frac{a^2\sqrt3}{4}\);
\item правильный шестиугольник со стороной $a$: \(\displaystyle S_{6}=\frac{3\sqrt3}{2}a^2\).
\end{itemize}
Шестиугольник можно разбить на \textbf{шесть равносторонних треугольников}, поэтому
\[
S_6=6\cdot\frac{a^2\sqrt3}{4}=\frac{3a^2\sqrt3}{2}.
\]

\ex{Пример 1. Правильная шестиугольная призма}
Сторона основания равна $1$, боковое ребро равно $\sqrt3$. В правильной призме боковое ребро является высотой. Следовательно,
\[
V=\frac{3\sqrt3}{2}\cdot 1^2\cdot\sqrt3
=\frac{9}{2}=4{,}5.
\]

\topic{Как меняется объём при изменении размеров}
Если \emph{все линейные размеры основания} изменены в $k$ раз, а высота --- в $m$ раз, то
\[
S_{\text{осн},2}=k^2S_{\text{осн},1}, \qquad
\boxed{\frac{V_2}{V_1}=k^2m}.
\]
Основания должны оставаться \textbf{подобными}. Для правильного треугольника и правильного шестиугольника достаточно изменить их сторону в $k$ раз.

\begin{center}
\begin{tabularx}{\linewidth}{@{}XcX@{}}
\toprule
Изменение & Коэффициент & Влияние на объём \\
\midrule
Сторона основания в $3$ раза больше & $k=3$ & объём в $9$ раз больше при прежней высоте \\
Высота в $2$ раза меньше & $m=\frac12$ & объём в $2$ раза меньше при прежнем основании \\
Оба изменения вместе & $k^2m=\frac92$ & объём в $4{,}5$ раза больше \\
\bottomrule
\end{tabularx}
\end{center}

\clearpage
\topic{Два способа решения задачи на изменение объёма}
\textbf{Условие.} Объём прямой призмы с равносторонним треугольником в основании равен $6$. Сторону основания увеличили в $3$ раза, а высоту уменьшили в $2$ раза. Найдите новый объём.

\ex{Способ 1. Деление формул (универсальный)}
Обозначим старые размеры через $a_1,h_1$, новые --- через $a_2,h_2$. По условию
\[
a_2=3a_1,\qquad h_1=2h_2,\qquad V_1=6.
\]
Запишем формулы для двух объёмов:
\[
V_1=\frac{\sqrt3}{4}a_1^2h_1,
\qquad
V_2=\frac{\sqrt3}{4}a_2^2h_2.
\]
Разделим второй объём на первый:
\begin{align*}
\frac{V_2}{V_1}
&=\frac{\frac{\sqrt3}{4}a_2^2h_2}
{\frac{\sqrt3}{4}a_1^2h_1}
=\frac{a_2^2h_2}{a_1^2h_1}\\
&=\frac{(3a_1)^2h_2}{a_1^2\cdot 2h_2}
=\frac92.
\end{align*}
Отсюда \(V_2=\frac92\cdot6=\boxed{27}\).

\ex{Способ 2. Экспресс-метод}
\begin{enumerate}
\item Сторона основания выросла в $3$ раза: площадь выросла в $3^2=9$ раз.
\item Высота уменьшилась в $2$ раза: коэффициент равен $\frac12$.
\item Перемножаем коэффициенты: $9\cdot\frac12=\frac92$.
\end{enumerate}
Итак, \(\displaystyle V_2=6\cdot9\cdot\frac12=\boxed{27}\).

\key{Проверка здравым смыслом:} увеличение площади основания в $9$ раз перевешивает уменьшение высоты вдвое, поэтому новый объём обязан быть больше исходного.

\topic{Дробь, составленная из двух дробей}
Если $b\ne0$, $c\ne0$, $d\ne0$, то
\[
\frac{\frac ab}{\frac cd}=\frac ab\cdot\frac dc=\frac{ad}{bc}.
\]
При делении дробей \textbf{вторую дробь переворачивают}. Этот приём позволяет упростить отношение двух формул объёма.

\clearpage
\topic{Когда подходит экспресс-метод}
Быстрый пересчёт коэффициентов особенно удобен, когда искомая величина выражена \textbf{произведением степеней параметров} и известен масштаб изменения каждого параметра.
\[
Q=C\,a^p h^q\quad\Longrightarrow\quad
\frac{Q_2}{Q_1}=k^p m^q,
\]
где коэффициент $C$ остаётся неизменным, $a_2=ka_1$, $h_2=mh_1$.

\key{Если формула содержит сумму,} сначала проверьте, как изменяются \emph{все} слагаемые. Например, периметр $a+b+c$ нельзя умножать на коэффициент изменения только одной стороны. Когда общий множитель для суммы выделить трудно, надёжнее записать две формулы и разделить их.

\ex{Пример 2. Боковая поверхность прямой призмы}
Для \textbf{прямой} призмы
\[
\boxed{S_{\text{бок}}=P_{\text{осн}}h}.
\]
При подобном изменении всего основания в $k$ раз периметр возрастает в $k$ раз, следовательно, при изменении высоты в $m$ раз
\[
\frac{S_{\text{бок},2}}{S_{\text{бок},1}}=km.
\]
Для объёма при тех же изменениях коэффициент будет $k^2m$. \textbf{Не путайте} зависимость от стороны в первой и второй степени.

\topic{Средняя линия треугольника}
Если $M$ и $K$ --- середины сторон $AB$ и $AC$ треугольника $ABC$, то
\[
AM=\frac{AB}{2},\qquad AK=\frac{AC}{2},\qquad
MK\parallel BC,\qquad MK=\frac{BC}{2}.
\]
Треугольники $AMK$ и $ABC$ подобны с коэффициентом $\frac12$, поэтому
\[
\boxed{\frac{S_{AMK}}{S_{ABC}}=\left(\frac12\right)^2=\frac14}.
\]
Эту закономерность используем при отсечении призмы плоскостью, которая проходит через среднюю линию основания \textbf{параллельно боковым рёбрам}.

\clearpage
\topic{Отсечение треугольной призмы}
Рассмотрим призму $ABCA_1B_1C_1$. Точки $M$ и $K$ --- середины $AB$ и $AC$. Через $MK$ проведена плоскость, параллельная боковым рёбрам. Она отделяет малую призму $AMKA_1M_1K_1$.

\begin{center}
\begin{tikzpicture}[scale=.97,line cap=round,line join=round]
\tkzDefPoint(0,0){A}
\tkzDefPoint(5.4,0){B}
\tkzDefPoint(1.45,2.15){C}
\tkzDefPoint(.78,2.55){Aa}
\tkzDefPoint(6.18,2.55){Ba}
\tkzDefPoint(2.23,4.70){Ca}
\tkzDefMidPoint(A,B)\tkzGetPoint{M}
\tkzDefMidPoint(A,C)\tkzGetPoint{K}
\tkzDefMidPoint(Aa,Ba)\tkzGetPoint{Ma}
\tkzDefMidPoint(Aa,Ca)\tkzGetPoint{Ka}
\fill[accent!10] (M)--(K)--(Ka)--(Ma)--cycle;
\tkzDrawSegments[black!60,line width=.75pt](A,B B,C C,A Aa,Ba Ba,Ca Ca,Aa A,Aa B,Ba C,Ca)
\tkzDrawSegments[accent,line width=1.4pt](M,K K,Ka Ka,Ma Ma,M)
\tkzDrawSegments[accent!80,line width=.9pt,dashed](A,M A,K Aa,Ma Aa,Ka)
\tkzLabelPoint[below left](A){$A$}
\tkzLabelPoint[below right](B){$B$}
\tkzLabelPoint[above left](C){$C$}
\tkzLabelPoint[left](Aa){$A_1$}
\tkzLabelPoint[right](Ba){$B_1$}
\tkzLabelPoint[above](Ca){$C_1$}
\tkzLabelPoint[below](M){$M$}
\tkzLabelPoint[left](K){$K$}
\tkzLabelPoint[below right](Ma){$M_1$}
\tkzLabelPoint[above left](Ka){$K_1$}
\end{tikzpicture}
\end{center}
\vspace{-2mm}
\textbf{Боковая поверхность: коэффициент $\frac12$.} Каждая боковая грань малой призмы имеет основание вдвое короче соответствующей грани большой призмы, а параллельное боковое ребро осталось тем же. Поэтому площадь каждой соответствующей боковой грани уменьшилась вдвое:
\[
\boxed{S_{\text{бок, мал}}=\tfrac12 S_{\text{бок, бол}}}.
\]
Свойство верно и для наклонной призмы: площади соответствующих параллелограммов масштабируются вдвое по длине основания при неизменном боковом ребре и угле между ними.

\textbf{Объём: коэффициент $\frac14$.} Основание малой призмы составляет четверть основания исходной, а перпендикулярная высота осталась прежней:
\[
\boxed{V_{\text{мал}}=\tfrac14 V_{\text{бол}}},\qquad
\boxed{V_{\text{оставшейся части}}=\tfrac34V_{\text{бол}}}.
\]

\clearpage
\topic{Два примера с отсечением призмы}
\ex{Пример 3. Найдите боковую поверхность меньшей призмы}
Боковая поверхность исходной треугольной призмы равна $36$. Через среднюю линию основания проведена плоскость, параллельная боковым рёбрам, которая отделяет меньшую призму.

\textbf{Решение через отношение:} соответствующие стороны оснований малой и большой призм относятся как $1:2$, а боковые рёбра равны и параллельны. Следовательно,
\[
\frac{S_{\text{бок, мал}}}{S_{\text{бок, бол}}}=\frac12,
\qquad S_{\text{бок, мал}}=\frac{36}{2}=\boxed{18}.
\]
\emph{Важно:} здесь отношение боковых площадей равно $\frac12$, а отношение объёмов --- $\frac14$. Причина --- боковая грань масштабируется по одному направлению, площадь основания --- по двум.

\ex{Пример 4. Найдите объём меньшей и оставшейся частей}
Объём исходной призмы равен $80$. Сечение устроено так же: через среднюю линию основания, параллельно боковым рёбрам.
\[
V_{\text{мал}}=\frac14\cdot80=\boxed{20},\qquad
V_{\text{ост}}=80-20=\boxed{60}.
\]
\textbf{Обобщение.} Если $AM/AB=AK/AC=t$ и $MK\parallel BC$ при $0<t<1$, а секущая плоскость параллельна боковым рёбрам, то
\[
\frac{S_{\text{бок, мал}}}{S_{\text{бок, бол}}}=t,
\qquad
\frac{V_{\text{мал}}}{V_{\text{бол}}}=t^2.
\]

\topic{Ошибки, на которые следует обращать внимание}
\begin{itemize}
\item \textbf{Перепутали масштаб стороны и площади.} При увеличении стороны основания втрое площадь подобного основания растёт в девять раз.
\item \textbf{Потеряли скобки.} $(3a)^2=9a^2$, а не $3a^2$.
\item \textbf{Перепутали высоту с ребром.} У наклонной призмы высота измеряется перпендикулярно основаниям.
\item \textbf{Без проверки применили $S_{\text{бок}}=Ph$.} Эта формула в таком виде относится к прямой призме.
\item \textbf{Подставили $\frac12$ вместо $\frac14$.} При отсечении по средней линии боковая поверхность и объём изменяются по разным законам.
\end{itemize}

\clearpage
\thispagestyle{fancy}
\begin{center}
{\Large\color{darkaccent}Алгоритм: как выбрать способ сравнения призм}\par
\vspace{9mm}
\begin{tikzpicture}[
font=\small, node distance=11mm,
start/.style={ellipse,draw=accent,very thick,align=center,minimum width=37mm,minimum height=10mm,inner sep=4pt},
act/.style={rectangle,rounded corners=2mm,draw=accent,thin,align=center,text width=95mm,minimum height=15mm,inner sep=5pt},
choice/.style={diamond,draw=accent,thin,aspect=3.0,align=center,text width=56mm,inner sep=2pt},
ar/.style={-Latex,thick,draw=darkaccent}]
\node[start] (start) {Начало};
\node[act,below=of start] (formula) {Запишите формулу искомой величины: $V$, $S_{\text{бок}}$ или $S_{\text{осн}}$};
\node[choice,below=14mm of formula] (decision) {Все множители можно пересчитать по известным коэффициентам?};
\node[act,text width=51mm] (fast) at ($(decision.south)+(-3.35cm,-3.0cm)$) {Экспресс-метод:\\ учтите степени и перемножьте коэффициенты};
\node[act,text width=51mm] (exact) at ($(decision.south)+(3.35cm,-3.0cm)$) {Метод отношений:\\ разделите формулы для двух фигур};
\node[act] (check) at ($(decision.south)+(0,-6.0cm)$) {Подставьте исходные данные и проверьте правдоподобие результата};
\node[start,below=of check] (end) {Ответ};
\draw[ar] (start)--(formula);
\draw[ar] (formula)--(decision);
\draw[ar] (decision.south west) |- node[pos=.18,above]{да} (fast.north);
\draw[ar] (decision.south east) |- node[pos=.18,above]{нет} (exact.north);
\draw[ar] (fast.south) |- ([xshift=-13mm,yshift=9mm]check.north) -- ([xshift=-13mm]check.north);
\draw[ar] (exact.south) |- ([xshift=13mm,yshift=9mm]check.north) -- ([xshift=13mm]check.north);
\draw[ar] (check)--(end);
\end{tikzpicture}
\end{center}
\clearpage

\section*{Тренировочный блок}
\noindent\textcolor{accent}{\rule{\linewidth}{0.45pt}}
\textbf{Путь А. 10 заданий от простого к более сложному.} В заданиях с сечениями меньшей называем призму, содержащую вершину $A$ исходного треугольного основания.

\subsection*{Базовый уровень}
\begin{enumerate}
\item Найдите площадь равностороннего треугольника со стороной $4$.
\item Найдите площадь правильного шестиугольника со стороной $2$.
\item Найдите объём правильной треугольной призмы, если сторона основания равна $2$, а высота равна $3\sqrt3$.
\item Найдите объём правильной шестиугольной призмы, если сторона основания равна $2$, а высота равна $\sqrt3$.
\end{enumerate}

\subsection*{Средний уровень}
\begin{enumerate}[start=5]
\item Объём прямой призмы равен $12$. Все линейные размеры её основания увеличили в $2$ раза, а высоту уменьшили в $2$ раза. Найдите новый объём.
\item Объём правильной треугольной призмы равен $8$. Сторону основания увеличили в $3$ раза, а высоту уменьшили в $3$ раза. Найдите новый объём.
\item После увеличения стороны правильной треугольной призмы в $2$ раза и уменьшения высоты в $3$ раза её объём стал равен $20$. Найдите первоначальный объём.
\end{enumerate}

\subsection*{Уровень выше среднего}
\begin{enumerate}[start=8]
\item Боковая поверхность треугольной призмы равна $50$. Через среднюю линию основания проведена плоскость, параллельная боковым рёбрам. Найдите боковую поверхность меньшей призмы.
\item Объём меньшей призмы, отсечённой от треугольной призмы плоскостью через среднюю линию основания параллельно боковым рёбрам, равен $14$. Найдите объём исходной призмы и оставшейся части.
\item В треугольной призме точки $M\in AB$ и $K\in AC$ выбраны так, что $AM/AB=AK/AC=\frac13$. Через $MK$ проведена плоскость, параллельная боковым рёбрам. Объём исходной призмы равен $81$, площадь её боковой поверхности равна $54$. Найдите объём и боковую поверхность меньшей призмы.
\end{enumerate}

\clearpage
\section*{Ответы к тренировочному блоку}
\noindent\textcolor{accent}{\rule{\linewidth}{0.45pt}}
\vspace{5mm}
\renewcommand{\arraystretch}{1.35}
\begin{center}
\begin{tabularx}{.92\linewidth}{@{}cX@{}}
\toprule
\textbf{№} & \textbf{Ответ}\\
\midrule
1 & $4\sqrt3$\\
2 & $6\sqrt3$\\
3 & $9$\\
4 & $18$\\
5 & $24$\\
6 & $24$\\
7 & $15$\\
8 & $25$\\
9 & исходная призма: $56$; оставшаяся часть: $42$\\
10 & объём меньшей призмы: $9$; боковая поверхность: $18$\\
\bottomrule
\end{tabularx}
\end{center}
\vspace{11mm}
\section*{Итоговая памятка}
\noindent\textcolor{accent}{\rule{\linewidth}{0.45pt}}
\begin{tabularx}{\linewidth}{@{}>{\raggedright\arraybackslash}p{.35\linewidth}X@{}}
\toprule
\textbf{Что известно} & \textbf{Что делать}\\
\midrule
Площадь основания и высота & $V=S_{\text{осн}}h$\\
Изменение подобного основания в $k$ раз & площадь умножается на $k^2$\\
Изменение высоты в $m$ раз & объём дополнительно умножается на $m$\\
Средняя линия треугольника & стороны меньшего основания --- $\frac12$, его площадь --- $\frac14$ исходной\\
Сечение призмы через среднюю линию & боковая поверхность малой призмы --- $\frac12$, объём --- $\frac14$ исходных\\
Сложная формула для отношения & записать обе формулы, разделить, сократить общие множители\\
\bottomrule
\end{tabularx}
\vfill
\begin{center}
{\small\color{lightgray}Лёвин Артём Александрович --- эксклюзивно для Ксении Клыковой}
\end{center}
\end{document}
